\emph{证明系统}提供了刻画语句间后承关系与相容性的纯句法方法。
\emph{自然演绎}就是这样一个推导系统。其中的一个\emph{推导}由一棵公式树构成。
推导最顶端的公式是\emph{假设}。
为使推导正确，所有其他公式都必须通过若干\emph{推理规则}之一获得正确辩护。
这些规则成对出现：每个联结词和量词都有一条引入规则和一条消去规则。
例如，如果一个公式~$!A$ 通过 $\Elim{\lif}$ 规则获得辩护，
那么它前面的公式（即\emph{前提}）必须是 $!B \lif !A$ 和 $!B$（对某个~$!B$）。
有些推理规则还允许将假设\emph{解除}。
例如，如果通过 $\Intro{\lif}$ 从 $!B$ 推出 $!A \lif !B$，
那么在导出前提~$!B$ 的子推导中，所有作为假设出现的~$!A$ 都可以被解除，
并被赋予一个标签，该标签亦记录于该推理处。

如果存在一个末公式为~$!A$ 且所有假设均被解除的推导，则我们称 $!A$ 是一个定理，
并记作~$\Proves !A$。
如果所有未解除的假设都属于某个集合~$\Gamma$，则我们称 $!A$ \emph{可从}~$\Gamma$ \emph{推导}，
并记作 $\Gamma \Proves !A$。
如果 $\Gamma \Proves \lfalse$，我们称 $\Gamma$ 是\emph{不一致的}，否则称其为\emph{一致的}。
这些概念互相关联，例如，$\Gamma \Proves !A$ 当且仅当 $\Gamma \cup \{\lnot !A\} \Proves
\lfalse$。
它们也与相应的语义概念相关，例如，如果 $\Gamma \Proves !A$ 那么 $\Gamma \Entails !A$。
自然演绎的这一性质——从~$\Gamma$ 可推导出的公式都保证为~$\Gamma$ 的语义后承——称为\emph{可靠性}。
\emph{可靠性定理}通过对推导长度进行归纳来证明，表明每次单独的推理都保持了语义后承关系：
只要各前提为各自未解除假设的语义后承，其结论亦为未解除假设的语义后承。